\SourceLeaf{250}%
\noindent vollständig in dreiseitige Pyramiden zerlegen, die in beiden Körpern paarweise symmetrisch sind.

II. Um jede dreiseitige Pyramide lässt sich eine Kugel konstruieren, deren 4 Radien sie mithilfe von Ebenen durch je zwei Radien in je 4 gleichseitige Pyramiden zerlegen. (Dabei können auch negative Pyramiden vorkommen.)

III. Führt man diese Konstruktion bei jedem Paar der unter I. erwähnten dreiseitigen Pyramiden aus, so entstehen lauter gleichseitige Pyramiden, die wieder paarweise symmetrisch sind.

IV. Jede gleichseitige Pyramide lässt sich mithilfe einer Senkrechten von der Spitze auf die Grundfläche und durch sie gelegter Ebenen in \textls{drei} dreiseitige Pyramiden zerlegen. Ihre Grundfläche ist ein gleichschenkliges Dreieck, und die eine Kante, die von der Spitze dieses gleichschenkligen Dreiecks ausgeht, steht senkrecht auf der Grundfläche.

V. Zwei Pyramiden der zuletzt beschriebenen Art, die in beiden Polyedern entsprechend liegen, sind nicht nur symmetrisch, sondern auch kongruent. Denn jede lässt sich durch eine Ebene in zwei zueinander symmetrische Hälften teilen.

VI. Demnach lassen sich beide Polyeder als (algebraische) Summen von lauter kongruenten Pyramiden auffassen und haben daher gleiches Volumen. Q.E.D.

Ich muss noch hinzufügen, dass der Beweis nicht ganz von mir stammt. Beim Nachdenken über Ihre Äußerung war ich so weit gekommen, dass ich die Teile III---V klar überblickte. Zunächst fiel es mir aber schwer, jede dreiseitige Pyramide auf gleichseitige zurückzuführen. Vielleicht hätte ich die Sache auf sich beruhen lassen, wenn sie nicht in einem Gespräch mit Dr.\ \textsc{Stegmann} zufällig zur Sprache gekommen wäre. Er hatte den glücklichen Einfall, die Kugel um die dreiseitige [Pyramide] zu Hilfe zu nehmen, womit der Schlüssel eigentlich gefunden war.

Soweit ich sehe, liegt in dem Vorstehenden kein Fehlschluss. Sie würden mich aber sehr verbinden, wenn Sie den Beweis genau prüften und mir gelegentlich Ihr Urteil darüber schrieben. \mbox{.\ .\ .\ .\ .\ .\ .}\}
